\chapter{Serlets y JSP}
\section*{3a-JSP}
\subsection*{Dudas de conocimiento}


Las ideas que hay que sacar de esta parte son:
\begin{itemize}
    \item  Def de Servlets
    \item  HttoServletRequest y HttpServletResponse
    \item variables implicitas 
    \item Scopes de un JSP
\end{itemize}


\begin{itemize}
    \item El concepto de JSP por lo que he entendido es simplemente un html con codigo que si se ha pedido antes se entrega directamente y si no se ha ejecutado antes, se compila y luego se envia el html. Segun el diagrama de la pagina 14 
\end{itemize}
\subsection*{Dudas de contenido para el examen}
\begin{enumerate}
    \item Hay que saberse la sintaxis de memoria de los metodos de HttpServletRequest y HttpServletResponse, y los servlets y elementos escripts o tan solo conocerlos y saber usarlos. 
    \item ¿Hay que saberse todos los metodos de los objetos implicitos de JSP? Durante las diapostivas hay varios signos de exclamación, entre ellos hay algunas diapositivas como la 39 que tiene exclamación y son metdos, esos son los metodos que nos debemos saber de memoria o nos debemos saber todos?
\end{enumerate}

\section*{3b-JSP}
\subsection*{Dudas de conocimiento}

\begin{itemize}
    \item cual es la diferencia exacta entre forward y redirect (pag 45)
    \item No entiendoa que se refiere en la diapo 51 de dierencias entre forward y redirect.
    \item ¿Cuando por forma general el navegador tiene que cambiar su url? para diferenciar en cuando usar redirect oforward (diapo 57)
    \item ¿hay que saberse, es decir, puede entrar las diapos del estilo de la diapo 66, es decir consejos de buen uso pero de cosas que no se suelen hacer?
\end{itemize}

